\begin{animation}[id="persistent-segtree-cses", label="CSES Range Queries and Copies -- Persistent Segment Tree on [1, 3, 2, 5]"]
\shape{arr}{Array}{size=4, data=[1,3,2,5], labels="0..3", label="Array (Version 0)"}
\shape{st}{Tree}{data=[1,3,2,5], kind="segtree", show_sum=true}
\shape{info}{Array}{size=1, data=[0], label="Current Version"}

\step
\narrate{Build the initial segment tree (Version 0) from array [1, 3, 2, 5]. The root stores sum=11. Internal nodes store partial sums: [0,1]=4, [2,3]=7. Leaves store individual elements.}

\step
\recolor{st.node["[0,3]"]}{state=done}
\recolor{st.node["[0,1]"]}{state=done}
\recolor{st.node["[2,3]"]}{state=done}
\narrate{Version 0 complete. The segment tree has 7 nodes: root [0,3] sum=11, left child [0,1] sum=4, right child [2,3] sum=7, and 4 leaves. All nodes shown in green (done) belong to V0.}

\step
\highlight{st.node["[0,3]"]}
\narrate{Query: sum(0, 3) on Version 0. Start at root [0,3]. The query range [0,3] fully covers node range [0,3], so return sum=11 immediately. No recursion needed.}

\step
\recolor{st.node["[0,3]"]}{state=idle}
\recolor{st.node["[0,1]"]}{state=idle}
\recolor{st.node["[2,3]"]}{state=idle}
\highlight{st.node["[0,3]"]}
\narrate{Query: sum(1, 2) on Version 0. Start at root [0,3]. Query range [1,2] partially overlaps, so we must recurse into both children.}

\step
\recolor{st.node["[0,3]"]}{state=dim}
\highlight{st.node["[0,1]"]}
\highlight{st.node["[2,3]"]}
\narrate{Visit [0,1] and [2,3]. Neither is fully inside [1,2], so recurse deeper into their children to find the exact leaves.}

\step
\recolor{st.node["[0,1]"]}{state=dim}
\recolor{st.node["[2,3]"]}{state=dim}
\recolor{st.node["[0,0]"]}{state=done}
\recolor{st.node["[1,1]"]}{state=good}
\recolor{st.node["[2,2]"]}{state=good}
\recolor{st.node["[3,3]"]}{state=done}
\narrate{Leaf [0,0] is outside query range, skip. Leaf [1,1]=3 is inside, take it. Leaf [2,2]=2 is inside, take it. Leaf [3,3] is outside, skip. Answer = 3 + 2 = 5.}

\step
\recolor{st.node["[0,3]"]}{state=idle}
\recolor{st.node["[0,1]"]}{state=idle}
\recolor{st.node["[2,3]"]}{state=idle}
\recolor{st.node["[0,0]"]}{state=idle}
\recolor{st.node["[1,1]"]}{state=idle}
\recolor{st.node["[2,2]"]}{state=idle}
\recolor{st.node["[3,3]"]}{state=idle}
\narrate{Now the key operation: Update a[1] = 7 on Version 0. In a persistent segment tree, we do NOT modify existing nodes. Instead, we create NEW nodes along the root-to-leaf path (path copying). Old nodes stay intact for V0.}

\step
\highlight{st.node["[1,1]"]}
\recolor{st.node["[1,1]"]}{state=current}
\narrate{Target: leaf [1,1]. The path from root to this leaf is: [0,3] -> [0,1] -> [1,1]. We will create 3 new nodes (one per level) for Version 1. The old nodes remain as V0.}

\step
\recolor{st.node["[0,3]"]}{state=current}
\recolor{st.node["[0,1]"]}{state=current}
\recolor{st.node["[1,1]"]}{state=current}
\recolor{st.node["[2,3]"]}{state=dim}
\recolor{st.node["[0,0]"]}{state=dim}
\recolor{st.node["[2,2]"]}{state=dim}
\recolor{st.node["[3,3]"]}{state=dim}
\narrate{Path copying in action: nodes marked current (blue) are the NEW nodes for V1. Nodes marked dim are SHARED with V0 -- no copying needed. New leaf [1,1] gets value 7. New [0,1] sum = 1+7 = 8. New [0,3] sum = 8+7 = 15.}

\step
\apply{arr.cell[1]}{value=7}
\recolor{arr.cell[1]}{state=current}
\recolor{st.node["[0,3]"]}{state=good}
\recolor{st.node["[0,1]"]}{state=good}
\recolor{st.node["[1,1]"]}{state=good}
\recolor{st.node["[2,3]"]}{state=dim}
\recolor{st.node["[0,0]"]}{state=dim}
\recolor{st.node["[2,2]"]}{state=dim}
\recolor{st.node["[3,3]"]}{state=dim}
\apply{info.cell[0]}{value=1}
\narrate{Version 1 created. The 3 new nodes (good/sky-blue) form V1's path: root sum=15, [0,1] sum=8, leaf [1,1]=7. The dim nodes [0,0]=1, [2,3]=7, [2,2]=2, [3,3]=5 are shared between V0 and V1. V0's old root (sum=11) still exists in memory. Only 3 new nodes = O(log n) space.}

\step
\recolor{st.node["[0,3]"]}{state=idle}
\recolor{st.node["[0,1]"]}{state=idle}
\recolor{st.node["[1,1]"]}{state=idle}
\recolor{st.node["[0,0]"]}{state=idle}
\recolor{st.node["[2,3]"]}{state=idle}
\recolor{st.node["[2,2]"]}{state=idle}
\recolor{st.node["[3,3]"]}{state=idle}
\recolor{arr.cell[1]}{state=idle}
\narrate{Copy operation: create Version 2 as a copy of Version 1. This is trivial -- just store roots[2] = roots[1]. Both versions share the ENTIRE tree. Cost: O(1) time, O(1) space. No nodes are allocated.}

\step
\apply{info.cell[0]}{value=2}
\recolor{st.node["[0,3]"]}{state=done}
\recolor{st.node["[0,1]"]}{state=done}
\recolor{st.node["[1,1]"]}{state=done}
\recolor{st.node["[0,0]"]}{state=done}
\recolor{st.node["[2,3]"]}{state=done}
\recolor{st.node["[2,2]"]}{state=done}
\recolor{st.node["[3,3]"]}{state=done}
\narrate{Version 2 is now active. All nodes are green (done) -- the entire tree is shared with V1. Array state for V2: [1, 7, 2, 5], sum=15. If we query V0, we still get the original sums (using V0's root).}

\step
\recolor{st.node["[0,3]"]}{state=idle}
\recolor{st.node["[0,1]"]}{state=idle}
\recolor{st.node["[1,1]"]}{state=idle}
\recolor{st.node["[0,0]"]}{state=idle}
\recolor{st.node["[2,3]"]}{state=idle}
\recolor{st.node["[2,2]"]}{state=idle}
\recolor{st.node["[3,3]"]}{state=idle}
\narrate{Now update a[3] = 10 on Version 2. Path from root to leaf [3,3]: [0,3] -> [2,3] -> [3,3]. Create 3 new nodes via path copying.}

\step
\recolor{st.node["[0,3]"]}{state=current}
\recolor{st.node["[2,3]"]}{state=current}
\recolor{st.node["[3,3]"]}{state=current}
\recolor{st.node["[0,1]"]}{state=dim}
\recolor{st.node["[1,1]"]}{state=dim}
\recolor{st.node["[0,0]"]}{state=dim}
\recolor{st.node["[2,2]"]}{state=dim}
\narrate{Path copying: 3 new nodes (current/blue) along root -> [2,3] -> [3,3]. Shared nodes (dim): [0,1], [0,0], [1,1], [2,2] -- these are V1's nodes reused by V2. New leaf [3,3] = 10, new [2,3] sum = 2+10 = 12, new root sum = 8+12 = 20.}

\step
\apply{arr.cell[3]}{value=10}
\recolor{arr.cell[3]}{state=current}
\recolor{st.node["[0,3]"]}{state=good}
\recolor{st.node["[2,3]"]}{state=good}
\recolor{st.node["[3,3]"]}{state=good}
\recolor{st.node["[0,1]"]}{state=done}
\recolor{st.node["[1,1]"]}{state=done}
\recolor{st.node["[0,0]"]}{state=done}
\recolor{st.node["[2,2]"]}{state=done}
\narrate{Version 2 updated. New nodes (sky-blue): root sum=20, [2,3] sum=12, [3,3]=10. Shared nodes (green): [0,1] sum=8, [0,0]=1, [1,1]=7, [2,2]=2. Total nodes across all versions: 7 (V0) + 3 (V1 update) + 3 (V2 update) = 13. Without persistence, 3 full trees = 21 nodes.}

\step
\highlight{st.node["[0,3]"]}
\narrate{Query: sum(0, 3) on Version 2. Start at V2's root (sum=20). Range [0,3] fully covers [0,3], return 20. Array for V2 is [1, 7, 2, 10].}

\step
\recolor{st.node["[0,3]"]}{state=good}
\recolor{st.node["[0,1]"]}{state=good}
\recolor{st.node["[2,3]"]}{state=good}
\recolor{st.node["[1,1]"]}{state=good}
\recolor{st.node["[0,0]"]}{state=good}
\recolor{st.node["[2,2]"]}{state=good}
\recolor{st.node["[3,3]"]}{state=good}
\recolor{arr.cell[3]}{state=good}
\recolor{arr.cell[1]}{state=good}
\narrate{Summary: 3 versions coexist. V0: [1,3,2,5] sum=11. V1: [1,7,2,5] sum=15. V2: [1,7,2,10] sum=20. Each update created only 3 new nodes (O(log 4) = O(2)). Copy was free. Total space: O(n + q*log n). Each query/update: O(log n). This is the persistent segment tree -- the standard technique for versioned array problems.}
\end{animation}
